\documentclass[11pt,a4paper]{article}
\usepackage[T1]{fontenc}
\usepackage{lmodern}
\usepackage[margin=24mm]{geometry}
\usepackage{amsmath}
\usepackage{booktabs,array,longtable}
\usepackage{xcolor}
\usepackage[hidelinks]{hyperref}
\usepackage{xurl}
\setlength{\parindent}{0pt}
\setlength{\parskip}{5pt}
\setlength{\emergencystretch}{2em}
\newcommand{\rev}[1]{\par\medskip\noindent\textit{\color{black!60}#1}\par\smallskip}
\newcommand{\resp}{\noindent\textbf{Response.} }

\begin{document}

{\Large\bfseries Response to the code and manuscript review of 23 October 2026\par}
\smallskip
{\large Declined, displaced, or just harder to catch? A simulation framework for attributing marine heatwave impacts on fished populations from catch per unit effort\par}
\smallskip
{\small Prepared 24 October 2026. Line numbers refer to the annotated PDF.\par}

\section*{Summary}

We thank the reviewer for a careful reading of both the code and the text, and for reproducing the full pipeline. Of the five code findings, one (finding 1) changed a computed quantity; we fixed it, added a test that would have caught it, and reran every stage. The other four are matters of statement, scope or housekeeping; for two of them we ran checks, now in the Supporting Information, that show the omitted effects are immaterial. Table~\ref{tab:moved} lists every reported number that moved after the rerun. The text annotations led to a global opening for the Introduction, a narrower statement of what ``fishery data'' means, a new paragraph on survey cost, several added references, and a large number of wording changes. The main text was shortened to meet the journal's word limit, with the material moved to a separate Supporting Information file.

\section*{Code findings}

\subsection*{1. The ``no redistribution'' state uses future heatwaves (High)}

\rev{With redistribution off, the fixed share is the mean of suitability over the whole simulated trajectory, which for the factual run includes events that have not yet occurred; the counterfactual averages a different trajectory, so the fixed shares differ between runs from month 1 and the paired contrast is non-zero before the first event.}

\resp Correct, and we are grateful for the catch. The Methods claimed climatological shares and the code did not deliver them. The fixed shares are now the mean suitability over the temperature series \emph{without} heatwaves, passed to both runs (\texttt{distribution\_shares()} gains \texttt{phi\_base} and \texttt{phi\_out\_base}; \texttt{run\_world()} gains \texttt{Tmat\_base} and \texttt{phi\_out\_base}; \texttt{simulate\_replicate()}, \texttt{exp\_temporal()} and the figure script pass them). A new test in \texttt{tests/testthat/test-core.R} checks, for three seeds of the decline scenario, that the two runs have identical initial shares and identical log abundance in every month before the first heatwave anomaly; all 62 tests pass.

Two points of interpretation. First, the reviewer's measure of the leak, the contrast before the first \emph{event}, counts something that is not a leak. The heatwave anomaly rises smoothly, and the stock is stressed by temperature itself, so an anomaly below the category one threshold that defines an event already changes abundance. Measured instead before the first \emph{anomaly}, the leak in the decline scenario had a median of 0.013 and a maximum of 0.032 over 40 replicates (0.001 to 0.002 in scenarios without decline); after the fix it is exactly zero in every replicate, whereas the contrast before the first event is unchanged because it is a real heatwave effect. We now say in Section~3.2.1 that sub-threshold anomalies stress the stock although they do not register in the severity index. Second, the effect on the reported results is small, because the leak was a slow drift that the trend term of the target regression largely absorbs. Table~\ref{tab:moved} gives every number that changed after rerunning all seven stages with the corrected code. No conclusion changes, but three statements were reworded to match the new numbers. With decline and fleet contraction both active, the nominal and delta indices now sit about 0.005 below the true effect (nominal $-0.028$ against a truth of $-0.024$), where before they were within 0.001; without contraction they remain within 0.001, and Section~4.2 now says so. The quarterly survey at 29 stations no longer matches the fishery indices (RMSE 0.091 against 0.058), so Section~4.4 no longer claims that a quarterly survey can match them at any station count tested. The correlation of the distributed lag fit with the true heatwave effect fell from 0.30 to 0.24, which strengthens rather than weakens the conclusion of Section~4.6 that the lag model describes without predicting better.

\subsection*{2. The delta indices do not standardise effort (Medium)}

\rev{Neither delta implementation includes effort, although the Methods describe prediction ``at a fixed reference effort''. The fitted month effects can absorb changing effort intensity, and either the estimators need an explicit effort treatment or the manuscript should describe them as effort-unadjusted.}

\resp Correct as a description, and we have taken the second option, with a check that the first would not have changed anything. The phrase ``at a fixed reference effort'' is removed; Section~3.5 now describes the delta index as a grid-cell by month model predicted over the full grid without an effort covariate, and says why that is adequate here: total effort is fixed, so effort per fished cell changes only when the footprint contracts, by 0.10 to 0.13 tows per cell during events from a mean of about 1.4 (about 8 percent). To confirm, we fitted an effort-adjusted delta index on 40 replicates of four scenarios with the manuscript seeds: a complementary log-log encounter model with $\log E$ as an offset, which is exact for the simulated encounter process, and $\log E$ as a covariate in the positive part, predicted at the mean effort per fished cell. The coefficients differ from the current index by at most 0.003, within the replicate standard error (Supporting Information, Section~S3 and Table~S3). The factorial analysis of variance says the same from the other side: fleet contraction, the only mechanism that changes effort allocation during events, has a main effect of $-0.001$ on the delta index. The check is reproducible with \texttt{scripts/review\_checks.R}.

\subsection*{3. Month 1 calculates fishing mortality but does not apply the state transition (Medium)}

\rev{Abundance in month 1 is fixed at $\log N_0$; the loop allocates effort and computes $F_1$ but applies the transition only for $t>1$, so there are 96 observations and 95 transitions. This needs to be stated, or it is an off-by-one error.}

\resp It is the initial-condition convention of a discrete-time state process, and it was not stated. Month 1 is the initial state; eq.~(9) generates months 2 to 96; effort is allocated and catch is observed in every month including the first. Both runs share the convention, so it cannot enter any paired contrast. Section~3.1 and the text after eq.~(9) now say so, and the unused month-1 mortality calculation has been removed from the loop so that the code matches the statement.

\subsection*{4. Thermal stress omits fish in the refuge}

\rev{Thermal stress sums over frame cells only; fish beyond the frame remain in total abundance but receive neither refuge thermal stress nor an explicit mortality process. This can change the total-versus-in-frame contrast and should be made explicit or reconsidered.}

\resp Made explicit, with a bound. The refuge lies one latitude band poleward of the frame and is cooler by the latitudinal gradient, so it rarely reaches the stress threshold. Over 40 replicates of the decline plus redistribution scenario, the refuge temperature exceeded $T_{\mathrm{crit}}$ in 0.08 percent of months (2.1 percent in the worst replicate); the stress that would have accumulated there was 0.02 percent of the total (0.7 percent at worst); and its cumulative effect on log abundance averaged 0.0006 (0.02 at worst), against heatwave effects of 0.01 to 0.03 per standard deviation of severity. The proportion of the stock outside the frame peaked at 13 percent on average. Section~3.2.3 now states that fish beyond the edge are neither fished nor stressed and gives the bound; Supporting Information Section~S2 gives the numbers; the Discussion notes that a warmer or equatorward refuge would make the assumption matter. On mortality: fishing mortality is removal based and proportional to the fleet's overlap with the stock inside the frame, so fish outside the frame are not caught, which is the intended behaviour; natural mortality is zero by default and would apply to the whole stock if set.

\subsection*{5. The cache is not keyed to run size, configuration or source version (Low)}

\rev{Cached stages are reused whenever the file exists, so a quick run after a full run reuses full stages while recording quick sizes, and a code change can leave stale results in place.}

\resp Fixed. \texttt{scripts/run\_all.R} now computes an md5 key from the run sizes, the configuration and every file in \texttt{R/}, stamps it on each cached stage, and recomputes any stage whose key differs. The key is written to \texttt{results/parts/cache\_key.txt}. The full rerun reported here was made with the new key from an empty cache.

\section*{Manuscript annotations}

\subsection*{Framing and scope}

\rev{The Western Australian example and crustacean fisheries come too early in the Introduction (lines 32 to 37); start global and set the scene for pelagic and other fisheries. The Fredston et al.\ result should frame the problem in the Introduction (line 543).}

\resp Done. The Introduction now opens with the global evidence on heatwave impacts and fisheries consequences, states the detection problem in survey data (Fredston et al., 2023) before turning to CPUE, and the Western Australian case returns to the Discussion as the application setting.

\rev{Catchability as usually defined already incorporates availability; give an equation and references (line 43).}

\resp Added. The Introduction now gives $\mathrm{CPUE} = qN$ with the standard definition of $q$ as a proportionality constant that absorbs availability as well as gear efficiency (Arregu\'in-S\'anchez, 1996; Wilberg et al., 2010), explains why we separate availability from gear efficiency (a heatwave acts on them differently), and gives eq.~(1) of the main text as the three-way product. Both references were checked against the originals.

\rev{Narrow the claims to the information actually represented: fisheries also provide size and age composition, tagging and controlled fishing experiments (line 71).}

\resp Agreed, and important. The Introduction now states that ``fishery data'' means the catch, effort and location records of the commercial fleet and that the identifiability result applies to that information set; the Discussion adds that composition, tagging and gear experiments carry information about catchability that catch rates do not, and are the natural way to relax the invariance without a survey. The abstract and conclusion say ``fishery catch and effort data''.

\rev{The survey costs money and many fisheries rely on commercial data alone (Conclusion).}

\resp A new paragraph in the Discussion (Implications for monitoring) meets this directly: the survey that outperformed every fishery index was small (a few stations visited monthly); what it buys is the ability to tell a decline from a catchability change at all; the relevant comparison is with the cost of the management error that fishery data alone cannot rule out; and where no survey is possible, the fishery-only remedies and their limits are restated.

\rev{A comment on local depletion in finfish versus crustaceans would be welcome (line 628).}

\resp Added to the limitations: the model excludes local depletion and therefore the sequential depletion of aggregations, a documented route to hyperstability in sedentary invertebrates and aggregating finfish, which would make the fishery signal still harder to read in those fisheries.

\rev{Remind the reader of heatwave duration before the survey design results (Figure 7). Is guidance on survey size relative to fishery and heatwave size given later (line 529)?}

\resp Section~4.4 now opens by noting that simulated events last one to four months, so a survey that visits less often can fall between events. Section~4.8 points back to Section~4.4 for the station counts, and Section~4.7 states that the survey's error scales with the size of the effect being estimated.

\rev{How susceptible is the boundary flow when background SST is already close to the thermal threshold (line 484)?}

\resp Two sentences added to Section~4.5: the direction depends on which side of the optimum the frame sits relative to the refuge; as a warming baseline carries the frame past the optimum, more of the stock leaves during events, which is what the sensitivity analysis shows (warming trend coefficient $-$0.24 on the survey's bias for total abundance).

\rev{Remove the companion study sentence for the journal (line 650).}

\resp Removed; it stays in the thesis chapter.

\subsection*{Presentation}

\rev{Figure 5 shows eight of the ten scenarios in Table 3; image text needs to be legible across all images; the grey dashed in-frame truth is hard to see; Figure 6 needs two sentences on interpretation; Figure 10 terms are not defined.}

\resp Figure~5 now shows all ten scenarios, with the spatial delta index added since it carries the redistribution result. Every figure is regenerated at a larger base font (12 pt) and full width. The in-frame truth is now an orange dashed line (the Okabe-Ito orange was free, so no colour carries two meanings). The Figure~6 caption gains two sentences on how to read the panels. Figure~10 (now Figure 8) uses the plain parameter names of Table~S1 and the caption points to it.

\rev{Correlations 0.09 and 0.29 are not in Table 5 (line 500).}

\resp They are correlations of the fitted series with the \emph{true} effect, which the table omitted. Table~S6 now has those columns beside the correlations with the fitted target.

\rev{Terms in eq.~(5) and elsewhere are defined after the equations or only in Table 1; K is not used in eqs.~5 to 8; units for eq.~(9); references for the thermal niche and movement assumptions; ``shares''; ``worlds''; ``processes''; ``roughly''; ``the fishery tracks''; ``Its''; ``hardest''; and the other wording items.}

\resp All applied. Every symbol in eq.~(5) is defined at the equation with a pointer to Table~1; $K$ is introduced where the episode schedule is described; the Table~1 caption gives units (rates per month, log abundance relative to carrying capacity, stress in degree-months); the ideal free distribution is cited for the movement rule (Fretwell and Lucas, 1969); ``shares'' is introduced as ``the proportion of the stock in each cell''; ``worlds'' is replaced by ``runs''; ``processes'' by ``components''; ``roughly'' removed and CPUE used throughout; ``fishing effort tracks''; ``hardest'' to ``most difficult''; the duplicated definition in eq.~(4) removed; the paragraph openings reworded. On ``Its'' (line 507): this is the possessive, which is correct, so it is unchanged. Table~S1 now has one decimal convention per row and a units column.

\rev{References to Supporting Information; check journal policy; the paper must stand alone.}

\resp Fish and Fisheries accepts Supporting Information, and we have used it to meet the word limit, but every number quoted in the main text now appears in a main-text table or figure as well; the supplement is cited only for full versions and for the checks described above.

\section*{Numbers that changed after the rerun}

\begin{table}[htbp]
\centering\small
\caption{Reported quantities before and after the correction of finding 1. All seven stages were rerun from an empty cache; the largest Monte Carlo standard error of a mean heatwave coefficient is 0.0025.}
\label{tab:moved}
\begin{tabular}{@{}lrrr@{}}
\toprule
Quantity & Before & After & Change \\
\midrule
True effect, decline (S1001) & $-$0.023 & $-$0.024 & $-$0.001 \\
Nominal, decline (S1001) & $-$0.022 & $-$0.028 & $-$0.005 \\
Delta, decline (S1001) & $-$0.024 & $-$0.029 & $-$0.004 \\
Survey, decline (S1001) & $-$0.024 & $-$0.024 & $-$0.001 \\
Nominal, catchability loss (S0011) & $-$0.038 & $-$0.038 & 0.000 \\
Delta, catchability loss (S0011) & $-$0.042 & $-$0.043 & $-$0.001 \\
Survey, catchability loss (S0011) & $-$0.000 & $-$0.001 & $-$0.001 \\
True effect, decline + catchability (S1011) & $-$0.023 & $-$0.023 & $-$0.001 \\
Nominal, decline + catchability (S1011) & $-$0.064 & $-$0.069 & $-$0.005 \\
Survey, decline + catchability (S1011) & $-$0.026 & $-$0.026 & $-$0.001 \\
In-frame effect, redistribution (S0101) & $-$0.010 & $-$0.010 & 0.000 \\
Survey, redistribution (S0101) & $-$0.009 & $-$0.009 & 0.000 \\
Spatial delta, redistribution (S0101) & $-$0.004 & $-$0.004 & 0.000 \\
Nominal, warm-edge loss (SPQ1) & $-$0.102 & $-$0.107 & $-$0.005 \\
Matched grid: nominal, loss only & $-$0.039 & $-$0.039 & 0.000 \\
Matched grid: nominal, mixture & $-$0.036 & $-$0.038 & $-$0.002 \\
Matched grid: nominal, decline only & $-$0.030 & $-$0.032 & $-$0.003 \\
Matched grid: true effect, decline only & $-$0.023 & $-$0.024 & $-$0.001 \\
Fishery CoG shift, redistribution (km) & 45.098 & 45.098 & 0.000 \\
Survey CoG shift, redistribution (km) & 48.929 & 48.929 & 0.000 \\
Fishery CoG shift, warm-edge loss (km) & 35.168 & 34.425 & $-$0.743 \\
Survey RMSE, 7 stations monthly & 0.042 & 0.043 & 0.001 \\
Survey RMSE, 14 stations monthly & 0.028 & 0.031 & 0.004 \\
Survey coverage, 14 stations monthly & 0.887 & 0.867 & $-$0.020 \\
Survey power, 14 stations monthly & 0.257 & 0.240 & $-$0.017 \\
Nominal RMSE, design scenario & 0.054 & 0.058 & 0.004 \\
Delta RMSE, design scenario & 0.060 & 0.065 & 0.004 \\
Survey RMSE, 29 stations quarterly & 0.059 & 0.091 & 0.032 \\
Distributed lag R2 (survey target) & 0.177 & 0.187 & 0.010 \\
Distributed lag corr. with truth & 0.295 & 0.241 & $-$0.054 \\
Distributed lag OOS RMSE & 0.495 & 0.545 & 0.051 \\
Same-month OOS RMSE & 0.405 & 0.437 & 0.032 \\
Key accuracy & 0.521 & 0.512 & $-$0.009 \\
Key indeterminate & 0.110 & 0.114 & 0.004 \\
Frame: in-frame effect, warming refuge & $-$0.010 & $-$0.010 & 0.000 \\
Frame: survey, static refuge & 0.008 & 0.008 & 0.000 \\
SRC nominal bias on catchability loss & $-$0.575 & $-$0.575 & 0.000 \\
Max MCSE & 0.003 & 0.003 & 0.000 \\
\bottomrule
\end{tabular}

\end{table}

\section*{References cited in this response}
Arregu\'in-S\'anchez, F. (1996). Catchability: a key parameter for fish stock assessment. \textit{Reviews in Fish Biology and Fisheries}, 6(2), 221--242.\\
Fredston, A. L. et al. (2023). Marine heatwaves are not a dominant driver of change in demersal fishes. \textit{Nature}, 621, 324--329.\\
Fretwell, S. D. and Lucas, H. L. (1969). On territorial behavior and other factors influencing habitat distribution in birds. \textit{Acta Biotheoretica}, 19, 16--36.\\
Wilberg, M. J., Thorson, J. T., Linton, B. C. and Berkson, J. (2010). Incorporating time-varying catchability into population dynamic stock assessment models. \textit{Reviews in Fisheries Science}, 18(1), 7--24.

\end{document}
